\documentclass[11pt,a4paper]{article}
\usepackage[margin=1.5cm]{geometry}
\usepackage{amsmath}

\setlength{\parindent}{0pt}
\pagenumbering{gobble}

\begin{document}

\begin{center}
    {\LARGE \textbf{Triangles Review --- Worked Solutions}}\\[4pt]
\end{center}

\begin{enumerate}

\item % Q1
By Pythagoras' theorem,
\[
x^2=4^2+3^2=16+9=25,
\]
so
\[
x=\sqrt{25}=5.
\]

\item % Q2
Angles in a triangle sum to \(180^\circ\):
\[
52^\circ+71^\circ+\theta=180^\circ
\]
\[
\theta=180^\circ-52^\circ-71^\circ=57^\circ.
\]

\item % Q3
In the right triangle,
\[
\tan 35^\circ=\frac{x}{7}.
\]
Hence
\[
x=7\tan 35^\circ\approx 4.90.
\]

\item % Q4
Using the cosine rule with the given labels:
\[
3^2=5^2+x^2-2(5)(x)\cos 40^\circ.
\]
So
\[
x^2-10x\cos 40^\circ+16=0.
\]
The discriminant is
\[
D=(10\cos 40^\circ)^2-64<0,
\]
so there is no real value of \(x\). The given data are inconsistent; no such triangle exists.

\item % Q5
The hypotenuse is \(7\), so
\[
x^2+3^2=7^2
\]
\[
x^2=49-9=40.
\]
Thus
\[
x=\sqrt{40}=2\sqrt{10}\approx 6.32.
\]

\item % Q6
By Pythagoras' theorem,
\[
x^2=5^2+2^2=25+4=29,
\]
so
\[
x=\sqrt{29}\approx 5.39.
\]

\item % Q7
Using the cosine rule,
\[
6^2=4^2+x^2-2(4)(x)\cos 55^\circ.
\]
So
\[
x^2-8x\cos 55^\circ-20=0.
\]
The positive solution is
\[
x=\frac{8\cos 55^\circ+\sqrt{64\cos^2 55^\circ+80}}{2}\approx 7.32.
\]

\item % Q8
The side \(x\) is opposite the \(40^\circ\) angle, so by the cosine rule,
\[
x^2=3^2+4^2-2(3)(4)\cos 40^\circ
\]
\[
x^2=25-24\cos 40^\circ.
\]
Therefore
\[
x=\sqrt{25-24\cos 40^\circ}\approx 2.57.
\]

\item % Q9
By the cosine rule,
\[
5^2=4^2+6^2-2(4)(6)\cos\theta.
\]
Thus
\[
25=52-48\cos\theta
\]
\[
48\cos\theta=27
\]
\[
\cos\theta=\frac{27}{48}=\frac{9}{16}.
\]
So
\[
\theta=\arccos\left(\frac{9}{16}\right)\approx 55.77^\circ.
\]

\item % Q10
Here \(AB=4\), \(AC=BC=3\), and \(\theta=\angle ACB\).
\[
4^2=3^2+3^2-2(3)(3)\cos\theta
\]
\[
16=18-18\cos\theta
\]
\[
18\cos\theta=2
\]
\[
\cos\theta=\frac{1}{9}.
\]
Therefore
\[
\theta=\arccos\left(\frac{1}{9}\right)\approx 83.62^\circ.
\]

\item % Q11
The angle \(\theta=\angle ABC\) is opposite the side of length \(4\), so
\[
4^2=3^2+6^2-2(3)(6)\cos\theta.
\]
Thus
\[
16=45-36\cos\theta
\]
\[
36\cos\theta=29
\]
\[
\cos\theta=\frac{29}{36}.
\]
So
\[
\theta=\arccos\left(\frac{29}{36}\right)\approx 36.34^\circ.
\]

\item % Q12
In the right triangle,
\[
\cos 32^\circ=\frac{6}{x},
\]
so
\[
x=\frac{6}{\cos 32^\circ}\approx 7.08.
\]

\item % Q13
Angles in a triangle sum to \(180^\circ\):
\[
40^\circ+65^\circ+\theta=180^\circ
\]
\[
\theta=75^\circ.
\]

\item % Q14
The third angle is
\[
180^\circ-43^\circ-34^\circ=103^\circ.
\]
By the sine rule,
\[
\frac{x}{\sin 43^\circ}=\frac{5}{\sin 34^\circ}.
\]
Hence
\[
x=\frac{5\sin 43^\circ}{\sin 34^\circ}\approx 6.10.
\]

\item % Q15
By Pythagoras' theorem,
\[
x^2=4^2+2.5^2=16+6.25=22.25=\frac{89}{4}.
\]
Thus
\[
x=\frac{\sqrt{89}}{2}\approx 4.72.
\]

\item % Q16
The angle \(\theta=\angle ACB\) is opposite \(AB=4\), with \(AC=7\) and \(BC=5\):
\[
4^2=7^2+5^2-2(7)(5)\cos\theta.
\]
So
\[
16=74-70\cos\theta
\]
\[
70\cos\theta=58
\]
\[
\cos\theta=\frac{58}{70}=\frac{29}{35}.
\]
Therefore
\[
\theta=\arccos\left(\frac{29}{35}\right)\approx 34.05^\circ.
\]

\item % Q17
In the right triangle, relative to \(\theta\), the opposite side is \(8\) and the adjacent side is \(6\):
\[
\tan\theta=\frac{8}{6}=\frac{4}{3}.
\]
So
\[
\theta=\arctan\left(\frac{4}{3}\right)\approx 53.13^\circ.
\]

\item % Q18
Using the cosine rule,
\[
6^2=5^2+x^2-2(5)(x)\cos 45^\circ.
\]
Since \(\cos 45^\circ=\frac{\sqrt2}{2}\),
\[
36=25+x^2-5\sqrt2\,x.
\]
Thus
\[
x^2-5\sqrt2\,x-11=0.
\]
Solving,
\[
x=\frac{5\sqrt2+\sqrt{94}}{2}\approx 8.38.
\]

\item % Q19
The side \(x\) is opposite the \(40^\circ\) angle, and the hypotenuse is \(10\):
\[
\sin 40^\circ=\frac{x}{10}.
\]
Hence
\[
x=10\sin 40^\circ\approx 6.43.
\]

\item % Q20
Angles in a triangle sum to \(180^\circ\):
\[
50^\circ+70^\circ+\theta=180^\circ
\]
\[
\theta=60^\circ.
\]

\end{enumerate}

\end{document}